\begin{examplebox}
	\textbf{\textcolor{CExample}{例1（基础）}}
	\textbf{\textcolor{CExample}{题目：}}
	在等比数列 $\{a_n\}$ 中，$a_1 = 2$，$q = 3$，求 $a_5$。
	\textbf{\textcolor{CExample}{解题步骤：}}
	\begin{description}[style=nextline, leftmargin=2.8em, labelsep=0.5em, itemsep=0.45em]
		\item[\textbf{第一步（代入公式）：}]
			由通项公式 $a_n = a_1 q^{\,n-1}$，令 $n=5$ 得
			\[
				\begin{aligned}
					a_5 &= 2 \times 3^{5-1} \\ &= 2 \times 3^4.
			\end{aligned}
			\]
			\par\textbf{理由：} 直接代入已知量 $a_1=2$, $q=3$, $n=5$。
		\item[\textbf{第二步（指数计算）：}]
			$3^4 = 81$，因此
			\[
				\begin{aligned}
					a_5 &= 2 \times 81 \\ &= 162.
			\end{aligned}
			\]
			\par\textbf{理由：} 指数运算。
		\item[\textbf{第三步（写出答案）：}]
			故 $a_5 = 162$。
			\par\textbf{理由：} 结论。
	\end{description}
	\textbf{\textcolor{CExample}{关键结论：}}
	\textbf{$a_5 = 162$}

	\medskip
	\textbf{\textcolor{CExample}{例2（稍变形）}}
	\textbf{\textcolor{CExample}{题目：}}
	在等比数列 $\{a_n\}$ 中，$a_3 = 12$，$q = 2$，求 $a_1$。
	\textbf{\textcolor{CExample}{解题步骤：}}
	\begin{description}[style=nextline, leftmargin=2.8em, labelsep=0.5em, itemsep=0.45em]
		\item[\textbf{第一步（变形公式）：}]
			由通项公式 $a_n = a_1 q^{\,n-1}$，取 $n=3$ 得
			\[
				a_3 = a_1 q^{2}.
			\]
			\par\textbf{理由：} 代入 $n=3$。
		\item[\textbf{第二步（列方程）：}]
			代入 $a_3=12$, $q=2$ 得
			\[
				\begin{aligned}
					12 &= a_1 \times 2^2 \\ &= 4 a_1.
			\end{aligned}
			\]
			\par\textbf{理由：} 将已知数值代入。
		\item[\textbf{第三步（解方程）：}]
			两边同除以 $4$，得
			\[
				\begin{aligned}
					a_1 &= \frac{12}{4} \\ &= 3.
			\end{aligned}
			\]
			\par\textbf{理由：} 系数化为 $1$。
	\end{description}
	\textbf{\textcolor{CExample}{关键结论：}}
	\textbf{$a_1 = 3$}

	\medskip
	\textbf{\textcolor{CExample}{例3（综合应用）}}
	\textbf{\textcolor{CExample}{题目：}}
	在等比数列 $\{a_n\}$ 中，$a_2 = 6$，$a_5 = 48$，求该数列的通项公式。
	\textbf{\textcolor{CExample}{解题步骤：}}
	\begin{description}[style=nextline, leftmargin=2.8em, labelsep=0.5em, itemsep=0.45em]
		\item[\textbf{第一步（列方程组）：}]
			由通项公式得
			\[
				\begin{cases}
					a_2 = a_1 q,\\ a_5 = a_1 q^{4}.
			\end{cases}
			\]
			\par\textbf{理由：} 分别代入 $n=2$ 和 $n=5$。
		\item[\textbf{第二步（消去 $a_1$）：}]
			两式相除，得
			\[
				\begin{aligned}
					\frac{a_5}{a_2} &= \frac{a_1 q^{4}}{a_1 q} \\ &= q^{3},
			\end{aligned}
			\]
			即 $48 / 6 = 8$，所以 $q^{3}=8$，解得 $q = 2$。
			\par\textbf{理由：} 用第二式除以第一式，消去 $a_1$。
		\item[\textbf{第三步（回代求 $a_1$）：}]
			将 $q=2$ 代入 $a_2 = a_1 q$，得 $6 = a_1 \times 2$，所以
			\[
				\begin{aligned}
					a_1 &= \frac{6}{2} \\ &= 3.
			\end{aligned}
			\]
			\par\textbf{理由：} 解一元一次方程。
		\item[\textbf{第四步（写通项公式）：}]
			将 $a_1=3$, $q=2$ 代入通项公式，得
			\[
				a_n = 3 \times 2^{\,n-1}.
			\]
			\par\textbf{理由：} 代入 $a_1$ 和 $q$。
	\end{description}
	\textbf{\textcolor{CExample}{关键结论：}}
	\textbf{$a_n = 3 \cdot 2^{\,n-1}$}
\end{examplebox}
